\section{问题重述}

光伏微网需要在负载需求、可再生出力和外部电网收费之间安排购电与储能。电池可以把低价时段购入或暂时富余的电量转移到后续时段，但其容量、功率和充放电效率限制了调节能力。题目要求在满足负载的基础上尽可能节省购电费用，因此应把真实结算账单作为比较策略的主要依据。

四个问题逐步改变可用信息和决策机会。问题一给出确定性输入，求单日日前购电及储能调度；问题二面对未知负载和光伏，需要先预测再制定日计划，缺口通过五倍价格紧急补购；问题三在规定预报时点允许付费修订常规购电，并根据实测储电量控制电池；问题四进一步引入价格不确定性，分别推广问题二的静态采购和问题三的滚动调度。

全文按问题顺序展开，每问依次讨论问题分析、模型建立和求解结果。以确定性LP和非对称结算下的场景采购构成统一母模型，再在滚动采购中加入基于执行规则回放的候选筛选；通过筛选时点消融选择最终配置，并报告实际费用与紧急补购的权衡。问题四进一步使用价格与净负荷的联合场景。每一问先解释新增机制，再给出变量、目标与约束，最后结合实际执行结果和对照实验判断模型效果。

\section{模型假设}

\begin{enumerate}
 \item 将附件中的功率解释为以标记时刻为右端点的10分钟平均功率，以$\Delta t=1/6$小时换算电量。这是时间离散的建模解释，不是题目明确给出的物理事实。
 \item 充电、放电效率分别为$\eta_c=\eta_d=0.9$，往返效率为0.81；忽略线路损耗和自放电。电池外侧充放电功率均不超过5000 kW，储电量保持在1200--10800 kWh。
 \item 不向外网售电，无法利用的富余电量可免费弃置。主结算口径中原计划全额付费且取消部分不退款；另设退款解释作敏感性对照。题目未给出的寿命折旧、碳成本或投资费用不进入主目标。
 \item 每10分钟可以观测当期负载、光伏、电价及储电量并完成平衡。紧急购电补足剩余缺口，不另设题目未给出的外网容量上限；未来实际值及尚未发布的预报不能提前用于决策。
 \item 问题一首末储电量均为6000 kWh。其他问题仅在1月1日初始化，之后实际状态逐日连续传递，不在每天零时或2月1日重新充满电池。
\end{enumerate}

\section{符号说明与数据准备}

\subsection{主要符号与统一单位}

购电、充放电和储电量均以kWh计，电价以元/kWh计。以$p$表示电价、$c$表示充电量，避免成本与充电变量混用。跨日问题增加日期下标$D$，不致混淆时省略。

\begin{table}[htbp]\centering\small
\caption{主要符号及含义}\label{tab:symbols}
\begin{tabularx}{\linewidth}{p{0.19\linewidth}Xl}\toprule
符号 & 含义 & 单位\\\midrule
$D,t,s$ & 日期、时段$t=0,\ldots,143$、场景索引 & ---\\
$L_{D,t},R_{D,t}$ & 实际负载与光伏功率 & kW\\
$n_{D,t},p_{D,t}$ & 实际净负荷电量、实际交易电价 & kWh、元/kWh\\
$g_{D,t},a_{D,t}$ & 零时计划、最终生效常规购电量 & kWh\\
$c_t,d_t,e_t$ & 电池外侧充电量、放电量、时段初储电量 & kWh\\
$h_t,w_t$ & 紧急购电量、弃置电量 & kWh\\
$\pi_s,Q_s$ & 场景概率、场景剩余时域紧急费用 & ---、元\\
$\rho_t,\lambda$ & 参考备用量、CVaR权重 & kWh、---\\
$\beta_T,k_\rho$ & 终端偏离惩罚系数、备用缩放系数 & 元/kWh、---\\
$E_{\min},E_{\max},U$ & 储电量下上限、单段充放电量上限$5000/6$ & kWh\\
\bottomrule\end{tabularx}
\end{table}

\subsection{统一物理时间与模板列序}

每日保留144个10分钟时段，物理区间及净负荷定义为
\begin{equation}
 I_t=[t/6,(t+1)/6),\qquad n_{D,t}=(L_{D,t}-R_{D,t})\Delta t.
 \label{eq:time_energy}
\end{equation}
源标记00:10对应物理00:00--00:10，最后一个源标记次日00:00对应当日23:50--24:00。模型按物理自然日结算，而附件5购电页保留原始首标签0:10--0:20及末标签次日0:00--0:10，按源样本列序写入，不根据显示标签再次平移结果。

\begin{table}[htbp]\centering\small
\caption{计算时间与原模板显示的对应}\label{tab:physical_mapping}
\begin{tabularx}{\linewidth}{lLLL}\toprule
索引 & 源采样右端点 & 模型物理区间 & 原模板显示区间\\\midrule
$t=0$ & 00:10 & 00:00--00:10 & 0:10--0:20\\
$t=60$ & 10:10 & 10:00--10:10 & 10:10--10:20\\
$t=143$ & 次日00:00 & 23:50--24:00 & 次日0:00--0:10\\
\bottomrule\end{tabularx}
\end{table}

144段映射冻结于\texttt{data/time\_mapping.json/csv}，同时记录各工作簿原标签及单元格位置。题目指定10:00--10:10取$t=60$；四小时区间累加连续24段；同日相邻正紧急量合并为一个区间。这些操作均调用同一物理索引，热力矩阵与能量图也沿用该索引。

\subsection{评价数据及图表口径}

问题一使用附件1。问题二至四的评价和填报覆盖2025年2月1日至12月31日334天，一月用于历史积累与有限候选校准。各问题的具体预测及可用信息分别在对应模型中说明。同一问题的比较策略从共同2月初态出发，各自连续执行；不同问题保留各自已冻结一月运行形成的初态。

预测器与实时执行器沿用原冻结配置；候选筛选的适用问题及启用时点依据开发性回测确定。最终配置确定后固定，本文所报告结果属于开发期闭环评价。第二问与第四问日前版保留未筛选母模型，第三问与第四问滚动版在零时及全部合法日内修订时点启用筛选。决策过程只使用当时可得信息，与选型参考过开发期费用是两个不同层面的要求，不能据前者宣称未接触数据上的独立测试。

实际总费用只计算计划、调整及紧急三项账单，辅助终端罚项和风险权重不加入报表。正文表格和附件5均从同一执行数组生成，保留全部指定日期和逐日数据。新增热力图展示时序分布或策略差异，瀑布图分解实际费用变化；分位阴影带表示场景或历史窗口间差异，不作为未经估计的置信区间。
